\chapter{Metas, programación cuadrática y no lineal}
\label{chap:optimizacion-avanzada}

Este capítulo reúne tres extensiones distintas de la programación lineal. Las
metas ordenan varios objetivos; la programación cuadrática añade curvatura al
objetivo manteniendo restricciones lineales; la programación no lineal admite
expresiones generales. Elegir una familia cambia tanto la formulación como la
garantía que puede ofrecer el solver. No basta con comparar el número que
aparece como objetivo: hay que revisar prioridades, factibilidad y exactitud.

\section{Programación por metas}
\subsection{Prioridades y desviaciones}
Una organización puede buscar primero el cumplimiento de una meta y después
optimizar otra magnitud, sin sacrificar lo ya conseguido. En la programación
lexicográfica se optimiza \(g_1(x)\), se fija su valor óptimo y se optimiza
\(g_2(x)\) dentro de ese conjunto; se continúa así en el orden declarado.
No se suman automáticamente objetivos con grandes pesos. El orden expresa
una decisión de política: una mejora de una prioridad posterior nunca compensa
un empeoramiento de una anterior.

Para una meta \(a^\top x=t\), las variables de desviación satisfacen
\[
 a^\top x+n-p=t,\qquad n,p\geq0.
\]
\(n\) mide el déficit y \(p\) el exceso. Minimizar \(n+p\) penaliza ambas
direcciones; minimizar solo \(n\) tolera el exceso. Los coeficientes y las
unidades deben indicar la penalización deseada. QSB no crea estas ecuaciones
por inferencia: el usuario declara variables, objetivos y restricciones.

\subsection{Ejemplo: proteger un compromiso de producción}
Hay capacidad para diez unidades entre los productos \(A\) y \(B\). Se desea
primero aproximar \(A\) a ocho unidades y después maximizar \(B\):
\[
\begin{aligned}
 &\text{Prioridad 1: }\min(n_1+p_1),\\
 &\text{Prioridad 2: }\max B,\\
 &A+B\leq10,\qquad A+n_1-p_1=8,\\
 &A,B,n_1,p_1\geq0.
\end{aligned}
\]
La prioridad no obliga por sí misma a satisfacer exactamente la meta si
resulta imposible; minimiza su desviación sujeta a las restricciones duras.
En este ejemplo sí es posible alcanzar desviación cero. Una vez fijada,
\(A=8\), por lo que el máximo posterior es \(B=2\).

El archivo \path{manual/examples/opt_goal.json} contiene todas las variables,
las dos filas y la lista ordenada \texttt{goals}; cada prioridad tiene nombre,
\texttt{sense} y \texttt{coefficients}. Las cotas y los tipos se declaran en
arreglos alineados con \texttt{variableNames}. Ejecute:
\begin{lstlisting}
swift run qsb validate-goal-json manual/examples/opt_goal.json
swift run qsb solve-goal-json manual/examples/opt_goal.json --backend native
\end{lstlisting}
El documento de salida conserva el modelo, los metadatos \texttt{backend} y
la solución.

En \texttt{goalOutcomes} se muestran prioridad, sentido y valor.
\begin{table}[htbp]
\centering
\caption{Resultado lexicográfico del caso de producción}
\begin{tabular}{lr}
\toprule
Magnitud & Valor \\
\midrule
Prioridad 1: desviación total & 0 \\
Prioridad 2: producción de \(B\) & 2 \\
\(A\) & 8 \\
\(n_1\), \(p_1\) & 0, 0 \\
\bottomrule
\end{tabular}
\end{table}
La capacidad queda activa y la meta se cumple. Si se invierte el orden, el
primer objetivo produce \(B=10\); con ese valor fijado, \(A=0\), \(n_1=8\)
y la desviación es ocho. Este contraste permite discutir quién decide la
precedencia y qué compromiso protege.

\subsection{Flujo en macOS y carácter del solver}
Abra el archivo normalizado en \texttt{Goal Programming}. En
\texttt{Definition}, revise variables y cotas, la lista de prioridades y las
restricciones lineales. Los controles de subir y bajar cambian el orden de
una prioridad. Introduzca la ecuación de desviación como una igualdad y
verifique el signo negativo de \(p_1\). Pulse \texttt{Validate} y
\texttt{Run}; la vista nativa de prioridades muestra los objetivos ordenados,
las cantidades finales y la información del backend. El JSON ofrece el
registro exportable completo. Las tablas admiten copiar y pegar; las columnas
de coeficientes deben coincidir con el orden actual de variables. Las tablas
grandes usan selección de fila y grupos de coeficientes.

El backend nativo resuelve una secuencia de LP o problemas enteros, añadiendo
una igualdad para conservar cada óptimo previo. El caso continuo aparece
como \texttt{exact}; este término caracteriza el método, sujeto a tolerancias
de punto flotante. Las metas enteras usan un recorrido educativo de escala
pequeña, identificado como \texttt{fixtureScale}, y dependen de las
limitaciones de ramificación y acotación. La relajación especial de pares de
desviación enteros solo se aplica cuando ecuaciones integrales justifican que
los residuos serán enteros; no elimina la integridad de las decisiones.

\texttt{validate} no resuelve. \texttt{external} ejecuta cada prioridad con
HiGHS si está instalado; QSB conserva el orden y el fijado lexicográfico.
Si falta el ejecutable, hay un error explícito. Una prioridad posterior
infactible, un objetivo no acotado o el agotamiento de límites requieren
revisar el diagnóstico, no interpretar resultados parciales como el plan final.

\section{Programación cuadrática}
\subsection{Convención del objetivo y supuestos}
QSB representa el objetivo como
\[
 f(x)=x^\top Qx+c^\top x,
\]
con restricciones lineales, cotas y tipos de variables. El JSON utiliza
\texttt{quadraticMatrix} para \(Q\) y \texttt{linearCoefficients} para \(c\).
La convención no contiene el factor \(1/2\). Para una matriz simétrica, el
coeficiente del producto \(x_ix_j\), \(i\ne j\), es \(2Q_{ij}\): un término
\(4XY\) se representa con \(Q_{12}=Q_{21}=2\). No coloque 4 en ambas
celdas, porque duplicaría el efecto.

Una matriz simétrica definida positiva produce un objetivo estrictamente
convexo; su mínimo sobre una región convexa, si existe, es único. Una matriz
definida negativa produce el caso equivalente de máximo estrictamente
cóncavo. Las condiciones de estacionariedad y factibilidad permiten entonces
identificar el óptimo global. Una matriz indefinida cambia el problema y no
queda cubierta por esa garantía.

\subsection{Ejemplo: ajuste con capacidad limitada}
Se desea aproximar dos cantidades a 2 y 3, respectivamente, con capacidad
total cuatro. Eliminando la constante 13 de la distancia cuadrática se obtiene
\[
\begin{aligned}
 \min\quad & X^2+Y^2-4X-6Y\\
 \text{sujeto a}\quad &X+Y\leq4,\qquad X,Y\geq0.
\end{aligned}
\]
La matriz es la identidad y el vector lineal es \((-4,-6)\). La eliminación
de la constante conserva las cantidades óptimas, pero cambia el valor
objetivo reportado: la distancia original es \(f(X,Y)+13\).

El modelo está en \path{manual/examples/opt_qp.json}. Ejecute:
\begin{lstlisting}
swift run qsb validate-qp-json manual/examples/opt_qp.json
swift run qsb solve-qp-json manual/examples/opt_qp.json --backend native
\end{lstlisting}
El óptimo es \(X=1.5\), \(Y=2.5\), \(f=-12.5\), con
\texttt{Capacity} en \texttt{activeConstraints}. La distancia cuadrática
original es 0.5. Un objetivo negativo es perfectamente válido: procede de la
constante omitida y no representa una distancia negativa.

La solución sin capacidad sería \((2,3)\), que no es factible. Con multiplicador
\(\lambda\geq0\), las ecuaciones de estacionariedad son
\[
 2X-4+\lambda=0,\qquad2Y-6+\lambda=0.
\]
La capacidad activa exige \(X+Y=4\), de donde \(\lambda=1\), \(X=1.5\) y
\(Y=2.5\). Las cotas inferiores no son activas. La convexidad estricta y las
condiciones anteriores certifican el mínimo global del modelo continuo.
QSB muestra las restricciones activas, pero esta derivación del multiplicador
es una comprobación académica independiente.

\subsection{Flujo en macOS y límites}
Abra el JSON en \texttt{Quadratic Programming}. \texttt{Definition} separa
variables y cotas, matriz cuadrática y restricciones lineales. Introduzca
los coeficientes lineales en las filas de variables y la identidad en la
matriz; confirme los nombres de sus dos ejes. Añada \texttt{Capacity} con
coeficientes \((1,1)\), relación \(\leq\) y término derecho 4. Valide y
pulse \texttt{Run}. La solución nativa presenta objetivo, variables, cotas,
restricciones activas y exactitud; el JSON permite revisar el documento.

\texttt{Copy Table} y \texttt{Paste Table} trasladan la matriz completa.
Reordenar variables conocidas conserva ambos ejes por nombre; los encabezados
pegados deben coincidir con el orden vigente. Para más de doce variables,
el editor muestra una fila y grupos de columnas, aunque el intercambio de
la tabla sigue siendo completo. Esta capacidad del editor no implica que el
solver nativo admita problemas grandes.

El nativo continuo enumera conjuntos activos y resuelve sistemas KKT para
mínimos estrictamente convexos o máximos estrictamente cóncavos. La etiqueta
es \texttt{exact}, con cálculos numéricos. El recorrido rechaza más de veinte
candidatos de desigualdad, contando las cotas que convierte a desigualdades;
el costo de enumerar conjuntos crece rápidamente. El nativo entero admite
solo maximización cóncava diagonal con variables no negativas, hasta cuatro
variables, y enumera puntos en cotas finitas derivadas. Se identifica como
\texttt{fixtureScale}; no es un solver MIQP general.

El backend externo HiGHS admite la traducción de los casos continuos
estrictamente convexos/cóncavos soportados. QSB adapta la convención mediante
una Hessiana \(H=2Q\), porque el solver externo utiliza
\(\tfrac12 x^\top Hx\). La adaptación también conserva sentido, cotas y
restricciones, y verifica objetivo y factibilidad al leer la solución.
HiGHS no resuelve MIQP en esta integración; las variables enteras cuadráticas
requieren el alcance nativo descrito o un motor futuro. El modo
\texttt{validate} no calcula un óptimo.

\section{Programación no lineal}
\subsection{Expresiones y dominio}
La programación no lineal (PNL) admite un objetivo \(f(x)\) y restricciones
\(g_i(x)\ \{\leq,=,\geq\}\ b_i\). QSB recibe expresiones textuales,
variables ordenadas y cotas anulables. El evaluador reconoce suma, resta,
producto, división, potencias, paréntesis, multiplicación implícita y
\texttt{exp}, \texttt{log}, \texttt{sqrt}, \texttt{sin}, \texttt{cos}.
Use producto explícito en documentos didácticos para facilitar la lectura.
Los nombres se reconocen sin distinguir mayúsculas de minúsculas; conviene
mantener una escritura uniforme.

La sintaxis válida no garantiza un dominio válido para todos los puntos:
\texttt{log(X)} requiere \(X>0\), \texttt{sqrt(X)} requiere \(X\geq0\) y una
división requiere denominador no nulo. Las cotas deben corresponder al caso
y proteger el dominio cuando sea posible. Los ángulos de seno y coseno se
expresan en radianes. El JSON guarda \texttt{objectiveExpression} y, en cada
restricción, \texttt{expression}, \texttt{relation} y \texttt{rhs}.

\subsection{Ejemplo: distancia con restricción}
El archivo \path{manual/examples/opt_nlp.json} plantea
\[
 \min\ (X-3)^2+(Y-2)^2,\qquad X+Y\leq4,\qquad0\leq X,Y\leq4.
\]
El objetivo es la distancia cuadrática al punto \((3,2)\), que viola la
capacidad. La proyección sobre \(X+Y=4\) es \((2.5,1.5)\) y el valor
matemático es 0.5. Aunque este caso podría formularse como QP, sirve para
comparar una entrada por expresión y un método aproximado con una solución
analítica conocida.
\begin{lstlisting}
swift run qsb validate-nlp-json manual/examples/opt_nlp.json
swift run qsb solve-nlp-json manual/examples/opt_nlp.json --backend native
\end{lstlisting}
La ejecución verificada devuelve aproximadamente
\(X=2.500000002\), \(Y=1.499999999\), objetivo 0.4999999995 y violación
máxima \(5.0\times10^{-10}\). La evaluación de \texttt{Capacity} es
aproximadamente 4.0000000005. El campo \texttt{iterations} registra 361 iteraciones acumuladas de la
búsqueda seleccionada; no cuenta todas las evaluaciones de la búsqueda desde
los otros puntos iniciales. Este valor describe esta ejecución, no un límite
universal ni una medida comparable con iteraciones simplex.

La pequeña violación no autoriza exceder físicamente una capacidad. Es un
residuo numérico de la búsqueda; para implementación real, compruebe los
valores con la tolerancia y las unidades del caso. La salida contiene
\texttt{constraintEvaluations}, \texttt{maximumViolation},
\texttt{iterations}, objetivo y valores de variables, además del modelo y
metadatos. Reportar el objetivo sin el residuo ocultaría información relevante.

\subsection{Uso en macOS y garantía aproximada}
Abra el JSON en \texttt{Nonlinear Programming}. La definición permite editar
el sentido, la expresión objetivo, variables con cotas y restricciones con
expresión, relación y término derecho. Introduzca
\texttt{(X-3)\^{}2+(Y-2)\^{}2}, límites 0 y 4, y \texttt{X+Y} con
\(\leq4\). Use \texttt{Validate} antes de \texttt{Run}. Una expresión inválida
se conserva para corregirla. Variables y restricciones admiten intercambio
con hojas de cálculo; cada pegado sustituye una tabla completa cuando es válido.
La solución presenta el candidato, límites, evaluaciones y violación máxima,
junto con el carácter aproximado y el JSON exportable.

El backend nativo usa diferenciación automática, búsqueda desde varios puntos
y penalizaciones progresivas de restricciones. Es determinista, se limita a
cuatro variables y se identifica como \texttt{approximate}. También contiene
una búsqueda de baja dimensión sobre variedades de igualdad para casos
pequeños. Estos mecanismos no prueban optimalidad global para funciones no
convexas. Un candidato factible puede ser local, y una búsqueda fallida no
constituye por sí sola una prueba matemática de infactibilidad.

Las desigualdades estrictas procedentes de importaciones se normalizan a
fronteras no estrictas y generan una advertencia estructurada. El campo
\texttt{normalizedStrictInequalities} registra esa transformación; debe
revisarse si la frontera tiene un significado físico excluido por el caso.
\texttt{validate} diagnostica sin resolver. \texttt{external} está
explícitamente no disponible para PNL; seleccione el nativo para los problemas
educativos soportados.

\section{Cómo comparar y revisar resultados}
Conserve cada entrada junto con su salida y el backend utilizado. En metas,
revise el orden y cada valor fijado. En QP, compruebe la convención de la matriz
y la curvatura admitida. En PNL, revise dominios, cotas, evaluaciones y
violación máxima antes de hablar de una solución aplicable. La validación del
archivo y la convergencia de una búsqueda son evidencias distintas.

Las tres familias tienen definición nativa y vista JSON; los campos pendientes
siguen editables cuando la validación falla. En las ejecuciones nativas de
macOS, la resolución se realiza en segundo plano y se impiden ejecuciones
duplicadas mientras hay una en curso. Cambiar el modelo puede dejar
obsoleto el resultado de ese cálculo. No existe actualmente un control de
cancelación para esos trabajadores. La presencia de etiquetas accesibles en
los controles y resultados no sustituye una auditoría completa de VoiceOver y
navegación por teclado, que sigue pendiente para estos editores.

\section{Ejercicios}
\begin{enumerate}
 \item Invierta las prioridades del ejemplo de metas. Compruebe \(B=10\),
 \(A=0\) y desviación total 8. Explique qué compromiso se ha cambiado.
 \item Reduzca la capacidad del ejemplo de metas a 6. Calcule la desviación
 mínima de la primera prioridad y justifique por qué la segunda no puede
 recuperar producción de \(B\) sin empeorar la primera.
 \item En el QP aumente la capacidad a 5. Prediga \((X,Y)=(2,3)\),
 objetivo \(-13\) y distancia original cero. Contraste con QSB.
 \item Introduzca un término \(4XY\) en un objetivo de prueba. Escriba
 la matriz simétrica correcta y examine su curvatura antes de elegir solver.
 \item Reformule el objetivo PNL del ejemplo en forma QP, conservando
 restricciones y cotas. Compare cantidades y residuos; explique la diferencia
 entre su objetivo con constante y el objetivo cuadrático sin ella.
 \item Proponga una función de una variable con dos mínimos locales.
 Explique qué evidencia adicional necesitaría para atribuir optimalidad global
 a un candidato obtenido por el nativo PNL.
\end{enumerate}
